\documentclass[11pt,a4paper]{article}
\usepackage[utf8]{inputenc}
\usepackage[french]{babel}
\usepackage[T1]{fontenc}
\usepackage{lmodern}
\usepackage[margin=1.5cm]{geometry}

\begin{document}

% En-tête
\begin{center}
    {\Huge \textbf{Léo Bernard}} \\[0.2cm]
    {\Large \textbf{Développeur Full-Stack Senior}} \\[0.1cm]
    léo.bernard@email.com \ | \ 06 86 93 79 80
\end{center}

\vspace{0.3cm}

% Résumé / Profil
\section*{Profil Professionnel}
\noindent
Ingénieur logiciel polyvalent possédant une maîtrise approfondie du cycle de développement complet, de l'analyse fonctionnelle des besoins à la mise en production continue (CI/CD). Je suis particulièrement à l'aise dans les environnements agiles et j'excelle dans la résolution de problèmes algorithmiques complexes. J'ai activement participé à la migration d'architectures monolithiques vers des écosystèmes microservices basés sur le cloud, maximisant ainsi les performances globales et la fiabilité de systèmes critiques.

\vspace{0.3cm}

% Compétences
\section*{Compétences Techniques}
\noindent
\textbf{Agile / Scrum} $\bullet$ \textbf{React} $\bullet$ \textbf{GraphQL} $\bullet$ \textbf{Node.js} $\bullet$ \textbf{Git} $\bullet$ \textbf{Python} $\bullet$ \textbf{Vue.js} $\bullet$ \textbf{Java}

\vspace{0.3cm}

% Expériences
\section*{Expériences Professionnelles}
\noindent \textbf{Développeur Full-Stack Senior / Confirmé} --- \textit{BlaBlaCar} \hfill \textbf{10/2025 à Aujourd'hui} \\
Conception et intégration d'une architecture orientée événements avec Apache Kafka pour synchroniser les données entre plusieurs microservices. Création d'une bibliothèque de composants UI réutilisables sous Storybook, accélérant le développement des nouvelles fonctionnalités front-end de 40\%. Encadrement technique d'une équipe de 4 développeurs juniors, incluant la réalisation de revues de code rigoureuses et l'animation d'ateliers de conception. Optimisation des requêtes SQL et migration vers PostgreSQL, améliorant les temps de réponse de l'API de manière significative et réduisant la charge serveur de 30\%. \\[0.3cm]
\noindent \textbf{Développeur Full-Stack Senior} --- \textit{Leboncoin} \hfill \textbf{12/2024 à 08/2025} \\
Migration de l'infrastructure d'hébergement vers AWS (EC2, S3, RDS), garantissant une scalabilité automatique lors des forts pics de trafic. Audit de sécurité du code legacy et correction de plusieurs vulnérabilités critiques (injections SQL, failles XSS) identifiées par les équipes de pentest. Développement d'une API REST robuste en Node.js capable d'encaisser plus de 10 000 requêtes par seconde lors des pics de charge saisonniers. Intégration de solutions de paiement tierces (Stripe, PayPal) via des webhooks sécurisés et asynchrones pour la plateforme e-commerce principale. \\[0.3cm]
\noindent \textbf{Stage — Assistant Développeur Full-Stack Senior} --- \textit{Criteo} \hfill \textbf{12/2022 à 12/2024} \\
Refonte intégrale de l'interface utilisateur en utilisant React et TypeScript, ce qui a conduit à une augmentation de 25\% de la rétention client. Mise en place de tests end-to-end automatisés avec Cypress et Jest, faisant passer la couverture de code globale de 45\% à plus de 85\%. Mise en place d'un système de cache distribué avec Redis, diminuant la latence des requêtes API récurrentes de plus de 60\%. \\[0.3cm]


\vspace{0.3cm}

% Formations
\section*{Formation \& Diplômes}
\noindent \textbf{Diplôme d'Ingénieur en Développement Web et Mobile} \hfill \textbf{2021 - 2024} \\
\textit{INSA Lyon} \\[0.3cm]
\noindent \textbf{Classes Préparatoires (CPGE) en Filière Scientifique} \hfill \textbf{2019 - 2021} \\
\textit{Lycée Scientifique} \\[0.3cm]


\end{document}